\section{Diffusion Model 回顾：从图像到生物}

\subsection{扩散模型的核心原理}

在深入蛋白质设计之前，Amini 首先回顾了 Diffusion Model 的核心思想。扩散模型的基本原理是\textbf{通过学习从噪声中恢复数据来实现生成}，包含两个互逆的过程：

\textbf{1. 前向加噪过程 (Forward Noising Process)}

从原始数据空间（如图像）出发，通过迭代地添加噪声来逐步破坏数据：

$$x_0 \xrightarrow{\text{noise}} x_1 \xrightarrow{\text{noise}} x_2 \xrightarrow{\text{noise}} \cdots \xrightarrow{\text{noise}} x_T \sim \mathcal{N}(0, I)$$

其中：
\begin{itemize}
    \item $x_0$：原始干净数据
    \item $x_t$：第 $t$ 步加噪后的数据
    \item $x_T$：经过 $T$ 步后接近纯随机噪声
    \item 这一过程\textbf{不需要训练}，只是简单地逐步添加高斯噪声
\end{itemize}

\textbf{2. 反向去噪过程 (Reverse Denoising Process)}

训练一个神经网络来学习反向过程，从噪声中逐步恢复出数据：

$$x_T \xrightarrow{\text{denoise}} x_{T-1} \xrightarrow{\text{denoise}} \cdots \xrightarrow{\text{denoise}} x_0$$

训练任务可以被形式化为：给定时间步 $t$ 的含噪数据 $x_t$，预测更少噪声的版本 $x_{t-1}$，或者等价地，预测残差噪声 $\epsilon$：

$$\mathcal{L} = \mathbb{E}_{x_0, \epsilon, t} \left[ \| \epsilon - \epsilon_\theta(x_t, t) \|^2 \right]$$

其中：
\begin{itemize}
    \item $\epsilon$：实际添加的噪声
    \item $\epsilon_\theta$：神经网络预测的噪声
    \item $t$：随机采样的时间步
\end{itemize}

\begin{importantbox}{为什么 Diffusion Model 适合蛋白质设计？}
Diffusion Model 在图像生成领域已展现出三大突破性优势：
\begin{enumerate}
    \item \textbf{生成质量}：样本的逼真度和保真度极高
    \item \textbf{多样性}：能覆盖数据分布的广泛区域，不会 mode collapse
    \item \textbf{可控性}：支持条件生成，可以指定生成目标
\end{enumerate}
这些特性恰好对应蛋白质设计的核心需求：生成的蛋白质需要生物学上可行（质量）、覆盖多样化功能（多样性）、并满足特定设计规格（可控性）。
\end{importantbox}

\subsection{从连续域到离散域的挑战}

标准的 Diffusion Model 处理的是\textbf{连续数据}（如图像像素值），前向过程通过添加连续的高斯噪声实现。然而，蛋白质序列是\textbf{离散数据}——每个位置上是 20 种氨基酸中的一种——不能简单地"添加高斯噪声"到一个字符上。

这引出了一个关键的技术挑战：\textbf{如何在离散数据上定义扩散过程？}

\subsection{本章小结}

Diffusion Model 的核心是学习从噪声到数据的逆向过程。其在图像领域的成功为蛋白质设计提供了有力的范式借鉴，但蛋白质序列的离散性质要求对标准连续扩散框架进行根本性的改造。


